\scanpage{14}
$\Pi$ $n$-\unsure{quasi|pseudo}-polyèdre comb. \uline{épinglé} universel

$G=\unsure{\mathscr G_n|\mathfrak g_n|\mathfrak G_n}$ son groupe d'autom, engendré par $(\underline\tau_i)_{0\leq i\leq n}$

$S$ schéma de base

$\boxed{C_0}={}$ catégorie \inserted{(groupoïde)} des réalisations géom non dég. projectives sur $S$
\margin{\uline{NB} $C$ catégorie rigide, dont les classes d'isom. sont en corr. 1-1 avec $\Gamma(S,\mathcal O_S)^{2n+1}$ via $\{\lambda_0,\alpha_0,\lambda_1,\alpha_1,\ldots,\lambda_{n-1},\alpha_{n-1},\lambda_n\}$}
$\simeq{}$ Catégorie des triples $(X,(\tau_*\unsure{\overset{\mathrm f}{=}|=}(\tau_i)_{0\leq i\leq n},s_0)$, $X$ fibré projectif de rg $n+1$, sur $S$, $\tau_i$ une pseudoréflexion stricte de $X$ (d'où centre $h_i$, cocentre $H_i$) $0\leq i\leq n$, $s_0$ section de $X$ sur $S$, avec conditions
\[
\left\{\begin{aligned}
1)\ &\deleted{\unsure{h_*}}(h_i)_{0\leq i\leq n}\ \deleted{\text{engendrent}}\ \inserted{\text{lin. indép.}}\ \text{en chaque fibre}\\
&(\Rightarrow C\text{ hyperplan de }X)\\
1')\ &(H_i)\ \text{lin indép en chaque fibre}\\
&(\Rightarrow c,\text{ section de }X\\
2)\ &\tau_i\tau_j=\tau_j\tau_i\quad\text{si }0\leq i,j\leq n,\ i,j\text{ non consécutifs}\\
&(\underset{\substack{\text{compte tenu}\\\text{de }1,\deleted{\unsure{1'|5}}}}{\Longleftrightarrow}\ h_i\subset H_j\ \deleted{\unsure{h_j}}\text{ si }0\leq i,j\leq n,\ i,j\text{ non conséc.})\\
3)\ &h_i\notin H_{i+1}\ \text{en tte fibre}\qquad 0\leq i\leq n-1\\
&(\Rightarrow\deleted{\unsure{\tau_i\tau_{i+1}\ne\mathrm{id}}}\ \tau_i,\tau_{i+1}\text{ ne comm. pas})\\
4)\ &s_0\subset H_i\quad(1\leq i\leq n)\\
5)\ &s_0\notin H_0\quad\text{en tte fibre}\quad\deleted{\unsure{\text{ou bien}}}
\end{aligned}\right.
\]
\margin{\uline{NB} 1') est conséquence des autres conditions 1) 2) 3) 4) 5) ; mais en l'absence d'un $s_0$ et 4) 5), il faut imposer 1') (ce qui rend 1) à 3) syst. autodual)}
% slanted
\margin{ces conditions ne faisant pas intervenir $s_0$, et \unsure{impliquent|impliquant} l'existence \unsure{locale} \unsure{unique} de $s_0$ satisfaisant 4$^\circ$)}
\margin{NB Il y a lieu d'introduire aussi $C$ cat. des couples $(E,\varphi_G^E)$ ou $(E,\tau_*)$ (\uline{sans} rig. $s_0$ ....)}

$\boxed{C''}={}$ cat des triples $(\deleted{\unsure{E}}V,\varphi_G^{\deleted{\unsure{E}}V},\deleted{\unsure{a_0}}L_0)$, $V$ fibré vectoriel de rg $n+1$ sur $S$,
\[
\varphi_G^V:G\longrightarrow\operatorname{Aut}(V)
\]
hom de grpes, $L_0\subset V$ sous-fibré inv, tels que
\[
\left\{\begin{aligned}
1)\ &\forall\,0\leq i\leq n,\ \inserted{\tau_i^V\overset{\mathrm{dfn}}{=}}\varphi_G^V(\underline\tau_i)\text{ est une ps. réflexion de }V,\\
&\text{stricte si }1\leq i\leq n,\text{ et associée à }L_0\text{ si }i=0\\
&\text{(pas néc. stricte si }i=0).\ \text{On en conclut}\\
&\deleted{\unsure{\ell}}\ \text{des }L_i\subset V,\ L'_i\subset V'\ \text{(ss-fibrés inv)}\quad 1\leq i\leq n\\
2)\ &V=\bigoplus L_i,\quad\deleted{\text{les sous-}}\bigl[\deleted{2')\ \text{des }(L'_i)_{1\leq i\leq n}\text{ indép.}}\\
&\text{en chaque fibre (i.e. }\bigcap_{1\leq i\leq n}H_i\ \text{\uline{droite}}\ \inserted{\Delta_d}\ \text{sur chaque fibre)}\bigr]\\
\deleted{3)}\ &\deleted{L_0\subset H_i\ \text{si }i\geq2}\qquad\deleted{\text{où on pose }H_i=(L'_i)^\perp}\\
&\deleted{\text{(\uline{NB} est conséquence de }1)\,2)\text{ si }\tau_0^V\text{ ps. réflex stricte)}}
\end{aligned}\right.
\]
\margin{\uline{NB} 1) 2) $\Rightarrow L_i\subset H_j$ si $1\leq i,j\leq n$ $i,j$ non cons.}
\margin{\uline{NB} 2') est conséquence des autres conditions y inclus 3), et 3) est une conséquence des autres y inclus}

% archive sheet number: 13

